\section{Co-directed response of nominal PZT droplets to the measured ZNMF jet}
The experiments were separate. Final SI PZT droplet and ZNMF gas velocities are
used without recalibration, rescaling, sign tuning or addition of the fields.
The ZNMF downstream direction is rotated upward: $u_g=-V$, $v_g=U$;
$x=x_{ref}-Y_j$, $y=X_j+y_{PZT,low}-X_{j,first}+\delta_y$.
The nominal axis is aligned to the documented empirical PZT reference
$x_{ref}=14.54665\,$mm without straightening measured wandering.
Because the PZT physical outlet is not established, the lowest PZT measurement
plane (0.0856 mm) and first supported gas plane
(0.3232 mm in the gas frame) are made coincident,
with configurable $\delta_y=0\,$mm. This does not align physical exits.
Seven supported native PZT positions across the plume initialise both measured
velocity components once; these are not individually tracked droplets or number
concentration measurements. No travel from the atomiser surface is inferred.

For nominal $d=10\,\mu$m, motivated by $5\,\mu$m aperture times 1.89 and
rounding the $9.45\,\mu$m estimate, the equations are
\begin{align}
\dot{\boldsymbol q}&=\boldsymbol v_d,\\
\dot{\boldsymbol v}_d&=-g(1-\rho_a/\rho_w)\boldsymbol e_y+
\frac{f(Re_p)}{\tau_p}(\boldsymbol u_g-\boldsymbol v_d),\\
Re_p&=\rho_a d\lVert\boldsymbol u_g-\boldsymbol v_d\rVert/\mu_a,
\qquad\tau_p=\rho_w d^2/(18\mu_a).
\end{align}
The diameter is a configurable model estimate, not measured sizing. Existing
20-degree-C properties are retained. The drag multiplier is unity for $Re_p\leq0.01$
and $1+0.15Re_p^{0.687}$ otherwise, with $Re_p<800$. Gravity and buoyancy
remain in the upward-positive physical frame. Gas interpolation requires four
finite paired vectors per source cell; boundaries and gaps are never extrapolated.
This prescribed-mean-gas approximation neglects plume-induced gas changes, jet
unsteadiness, evaporation and droplet-size evolution. The combined flow is not
experimentally validated.

All 7 integrations complete. Initial sampled gas axial velocities
span -0.68911 to -0.19502 m/s;
the saved near-entry signed reversal region is retained, despite the nominal upward axis.
7 trajectories end with downward droplet velocity. Maximum heights
span 0.17519--0.334543 mm.
Cases exhibiting upward drag: 0; cases exhibiting
net upward axial acceleration: 0.
Over the common first-ascent interval $y=0.0856$--$0.17519\,$mm,
travel times span 0.0649192--0.2775 ms and axial velocity changes
span -0.771021---0.407629 m/s. These are common-interval comparisons, not
performance comparisons of differently terminated endpoints. The full histories
retain subsequent descent; branch styles distinguish ascent and descent.

Upward drag requires positive axial slip; net upward acceleration additionally
requires $a_{D,y}>g(1-\rho_a/\rho_w)$. Force-sign intervals are tabulated.
Net upward acceleration of a descending particle can mean deceleration of its
downward motion, not upward transport. The measured fixed-line PZT-only mean
profile is shown separately in red and is not a tracked single-droplet trajectory.

$\tau_p=0.305825\,$ms is a Stokes-limit time; the exported
$\tau_p/f(Re_p)$ is a finite-Re drag/slip coefficient timescale, not an exact
relaxation time. Scalars are tabulated. $Stk_f=\tau_p f=0.036699$ uses
ZNMF actuation $f=120\,$Hz without $2\pi$; neither acquisition
nor atomiser frequency is substituted. Local trajectory $Stk_D=\tau_p\lVert\boldsymbol u_g\rVert/D_0$
uses $D_0/\lVert\boldsymbol u_g\rVert$. Gas local $Re_{D,c}=U_cD_0/\nu_a$
and $Sr_{D,c}=fD_0/U_c$ use positive finite measured $U_c$ with retained gaps,
$D_0=2\,$mm, and do not imply nozzle-exit $Re_0$ or $Sr_0$. Stk and Sr are
distinct. Supporting PZT-only Re is conditional on explicitly assumed zero
background gas, not a measured ambient-gas condition. No collection efficiency,
lung deposition or combined-experiment validation is claimed.
